\chapter{Analyse et Spécification des Besoins}

\section*{Introduction}

Ce chapitre constitue le fondement conceptuel de la plateforme TalentFlow. Nous
identifions les acteurs du système et spécifions les besoins fonctionnels et non
fonctionnels. Nous présentons ensuite le diagramme de cas d'utilisation global, le
backlog du produit, le diagramme de classes et l'environnement de travail. Enfin,
nous détaillons l'architecture globale de la plateforme et le sprint backlog général.

% ============================================================
\section{Identification des Acteurs de Notre Application}
% ============================================================

La plateforme TalentFlow implique cinq catégories d'acteurs humains et un acteur système.
Le tableau~\ref{tab:acteurs} présente leurs responsabilités principales.

\begin{table}[h]
\centering
\renewcommand{\arraystretch}{1.6}
\begin{tabular}{|p{3cm}|p{3.5cm}|p{6.5cm}|}
\hline
\textbf{Acteur} & \textbf{Nature} & \textbf{Responsabilités principales} \\
\hline
\textbf{Guest} & Utilisateur non authentifié & Accès public, consultation des offres, soumission d'une candidature avec upload de CV. \\
\hline
\textbf{Candidat} & Utilisateur authentifié & Gestion du profil (\texttt{CandidateProfile}), sauvegarde d'offres (\texttt{SavedJob}), suivi des candidatures, passage des quiz et entretiens virtuels. \\
\hline
\textbf{Membre (Expert)} & Collaborateur RH du tenant & Consultation des candidatures assignées (\texttt{AssignationExpert}), évaluation (\texttt{AvisExpert}), consultation des résultats quiz et entretiens. \\
\hline
\textbf{Tenant (Recruteur)} & Entreprise cliente & Gestion des offres, candidatures, membres et entreprises ; planification des quiz et entretiens ; consultation des rapports IA. \\
\hline
\textbf{Admin} & Acteur technique interne & Validation et supervision des tenants, tableau de bord global, gestion du profil administrateur. \\
\hline
\textbf{AI Service} & Acteur système (Python/FastAPI) & Extraction PDF (PyMuPDF), OCR Tesseract, NER RoBERTa, embeddings et similarité sémantique (Sentence Transformers), matching et classement (ChromaDB, FAISS), résumé et raisonnement LLM (Groq), scoring composite, génération de questions, évaluation sémantique, prédiction de fraude et rapports. \\
\hline
\end{tabular}
\caption{Identification des acteurs de TalentFlow}
\label{tab:acteurs}
\end{table}

% ============================================================
\section{Étude des Besoins}
% ============================================================

\subsection{Les Besoins Fonctionnels}

\subsubsection{Besoins Fonctionnels du Guest}

\begin{itemize}
    \item \textbf{Accès à la plateforme :} naviguer sur l'interface publique sans compte.
    \item \textbf{Consultation des offres :} consulter les offres publiées (\texttt{EstPublie = true}).
    \item \textbf{Soumission d'une candidature :} postuler via le formulaire personnalisé
    (\texttt{FormulaireCandidature}) et uploader son CV sur Cloudinary (\texttt{CvUrl}).
\end{itemize}

\subsubsection{Besoins Fonctionnels du Candidat}

\begin{itemize}
    \item \textbf{Gestion du profil :} créer et mettre à jour son \texttt{CandidateProfile}
    (compétences, expériences, formations, certifications, langues, CV, liens, avatar).

    \item \textbf{Sauvegarde des offres :} sauvegarder des offres en favoris
    (\texttt{SavedJob}) et recevoir des notifications (\texttt{CandidatNotification}).

    \item \textbf{Suivi des candidatures :} consulter l'état de ses candidatures
    (\texttt{Candidature.Statut}) depuis son tableau de bord.

    \item \textbf{Passage du quiz technique :} accéder au quiz via \texttt{QuizToken},
    répondre aux QCM dans le temps imparti (\texttt{TimePerQuestion}) et consulter
    ses résultats (\texttt{QuizResult}).

    \item \textbf{Passage de l'entretien virtuel :} confirmer un créneau
    (\texttt{CreneauxDisponibles}), accéder via \texttt{LienToken} et interagir
    oralement avec l'avatar recruteur.
\end{itemize}

\subsubsection{Besoins Fonctionnels du Membre (Expert)}

\begin{itemize}
    \item \textbf{Consultation des offres assignées :} accéder aux offres via
    \texttt{AssignationExpert}.

    \item \textbf{Évaluation des candidatures :} consulter les CV, scores IA et
    résumés LLM dans \texttt{AnalyseCV} ; attribuer un score et rédiger un avis
    (\texttt{AvisExpert}).

    \item \textbf{Changement de statut :} modifier le statut d'une candidature.

    \item \textbf{Consultation des résultats :} accéder aux \texttt{QuizResult} et
    aux rapports \texttt{EntretienIA} incluant les indicateurs de fraude
    (\texttt{AlertesSecurite}, \texttt{NbChangementsOnglet}, \texttt{VerificationFacialeOk}).
\end{itemize}

\subsubsection{Besoins Fonctionnels du Tenant (Recruteur)}

\begin{itemize}
    \item \textbf{Gestion des entreprises :} créer et gérer les entités
    \texttt{Entreprise} (nom, secteur, RNE, logo).

    \item \textbf{Gestion du profil tenant :} renseigner le \texttt{TenantProfile}
    (titre, site web, réseaux sociaux, stack technique, statut de recrutement).

    \item \textbf{Gestion des membres :} inviter des \texttt{Expert}, les assigner
    à des offres (\texttt{AssignationExpert}).

    \item \textbf{Gestion des offres :} créer et publier des \texttt{OffreEmploi} ;
    personnaliser le formulaire de candidature (\texttt{FormulaireCandidature} /
    \texttt{ChampPersonnalise} : Texte, Nombre, Date, ChoixMultiple, Fichier) ;
    générer un lien public tokenisé (\texttt{PublicToken}).

    \item \textbf{Gestion des candidatures :} suivre le pipeline, accéder aux
    \texttt{AnalyseCV}, consulter les scores IA (matching sémantique, compétences,
    expérience, raisonnement LLM) et classer les candidats par offre, changer les statuts.

    \item \textbf{Planification du quiz :} paramétrer la date (\texttt{ScheduledDate}),
    le minuteur (\texttt{TimePerQuestion}) et les consignes ; consulter les
    \texttt{QuizResult}.

    \item \textbf{Planification de l'entretien :} définir les créneaux, activer le
    \texttt{LienToken} ; consulter le \texttt{RapportIA} avec indicateurs de fraude.

    \item \textbf{Rapports et exports :} générer des rapports détaillés ; exporter
    en PDF, Excel et Word.
\end{itemize}

\subsubsection{Besoins Fonctionnels de l'Admin}

\begin{itemize}
    \item \textbf{Gestion des tenants :} valider, approuver, rejeter ou suspendre les comptes
    tenants (\texttt{Tenant.Statut} : Pending, Approved, Rejected).
    \item \textbf{Administration de la plateforme :} consulter le tableau de bord global,
    gérer son profil administrateur et modifier son mot de passe.
\end{itemize}

\subsection{Les Besoins Non Fonctionnels}

\begin{table}[h]
\centering
\renewcommand{\arraystretch}{1.6}
\begin{tabular}{|p{3cm}|p{10cm}|}
\hline
\textbf{Critère} & \textbf{Exigence} \\
\hline
\textbf{Performance} & Réponse APIs $<$ 200ms. Analyse CV complète $<$ 30s.
Rendu 3D avatar à 60 FPS. \\
\hline
\textbf{Sécurité} & Authentification JWT. Chiffrement AES-256 (repos) et TLS 1.3 (transit).
Isolation logique par tenant. Accès quiz/entretien par tokens opaques.
Détection de fraude comportementale via TensorFlow.js et vérification faciale
(\texttt{AlertesSecurite}, \texttt{NbChangementsOnglet}, \texttt{VerificationFacialeOk}). \\
\hline
\textbf{Scalabilité} & Architecture modulaire. Conteneurisation Docker pour la montée
en charge horizontale. \\
\hline
\textbf{Disponibilité} & SLA 99,9\%. Reprise sur erreur pour le service IA Python/FastAPI. \\
\hline
\textbf{Maintenabilité} & Séparation claire des responsabilités. Documentation complète
des APIs. Pipeline CI/CD via GitHub Actions. \\
\hline
\textbf{Accessibilité} & Interface responsive desktop et mobile. Support FR et EN. \\
\hline
\textbf{Conformité RGPD} & Consentement avant enregistrement. Droit à l'effacement.
Minimisation des données. \\
\hline
\end{tabular}
\caption{Besoins non fonctionnels de TalentFlow}
\label{tab:nfr}
\end{table}

% ============================================================
\section{Diagramme de Cas d'Utilisation Global}
% ============================================================

Le diagramme de cas d'utilisation global (figure~\ref{fig:uc_global}) illustre les
interactions entre les cinq acteurs humains et l'acteur système AI Service. On y distingue :

\begin{itemize}
    \item Les fonctionnalités du Guest : accès plateforme, consultation des offres,
    soumission d'une candidature (\textit{Apply for an offer}, incluant \textit{Upload CV}).

    \item Les fonctionnalités du Candidat authentifié : gestion du profil, sauvegarde
    d'offres, suivi des candidatures, passage du quiz et de l'entretien virtuel.

    \item Les fonctionnalités du Membre : consultation des offres et candidatures assignées,
    scoring (\textit{Score candidate}), changement de statut, consultation des rapports.

    \item Les fonctionnalités du Tenant : gestion des offres (\textit{Manage offers}) avec
    personnalisation du formulaire et assignation de membres, planification des quiz et
    entretiens, gestion des candidatures.

    \item Les fonctionnalités de l'Admin : gestion des tenants et supervision de la plateforme.

    \item Les services de l'AI Service : extraction du texte des CV (PyMuPDF, complément
    OCR Tesseract si pages scannées), extraction structurée NER, similarité sémantique
    (Sentence Transformers), matching et classement (ChromaDB, FAISS), résumé et scoring
    via LLM (Groq), génération des questions de quiz et d'entretien, évaluation des
    réponses, prédiction de fraude et génération des rapports.
\end{itemize}

\begin{figure}[ht]
    \centering
    \includegraphics[width=0.85\textwidth]{images/usecase_diagram.png}
    \caption{Diagramme de cas d'utilisation global de TalentFlow}
    \label{fig:uc_global}
\end{figure}

% ============================================================
\section{Backlog du Produit}
% ============================================================

Le product backlog regroupe l'ensemble des fonctionnalités attendues de TalentFlow,
exprimées sous forme de user stories et priorisées selon leur valeur métier. Le
tableau~\ref{tab:backlog} présente le backlog réparti sur 8 sprints (Sprint 0 à Sprint 7).

\begin{table}[h]
\centering
\renewcommand{\arraystretch}{1.4}
\begin{tabular}{|c|p{8.5cm}|c|c|}
\hline
\textbf{ID} & \textbf{User Story} & \textbf{Priorité} & \textbf{Sprint} \\
\hline
US00 & En tant qu'équipe, nous voulons réaliser l'apprentissage technologique, concevoir
le design system, les maquettes UI/UX, la modélisation UML et le schéma de base de données
PostgreSQL. & Haute & 0 \\
\hline
US01 & En tant qu'administrateur, je veux mettre en place l'infrastructure multi-tenant
avec isolation logique des données par tenant. & Haute & 1 \\
\hline
US02 & En tant qu'administrateur, je veux gérer les comptes utilisateurs et les rôles
(Admin, Tenant, Expert, Candidat). & Haute & 1 \\
\hline
US03 & En tant qu'administrateur, je veux valider et configurer les tenants pour assurer
leur onboarding sur la plateforme. & Haute & 2 \\
\hline
US04 & En tant que tenant, je veux gérer les experts (\texttt{Expert}) et les entreprises
(\texttt{Entreprise}) de mon espace. & Haute & 2 \\
\hline
US05 & En tant que tenant, je veux créer et publier une offre d'emploi (\texttt{OffreEmploi})
afin d'attirer des candidats qualifiés. & Haute & 3 \\
\hline
US06 & En tant que tenant, je veux personnaliser le formulaire de candidature
(\texttt{FormulaireCandidature} / \texttt{ChampPersonnalise}) et générer un lien
public tokenisé (\texttt{PublicToken}). & Haute & 3 \\
\hline
US07 & En tant que guest ou candidat, je veux accéder à l'interface publique des offres
et soumettre ma candidature avec mon CV (hébergé sur Cloudinary). & Haute & 4 \\
\hline
US08 & En tant que tenant ou membre, je veux suivre le pipeline des candidatures reçues
et consulter les profils soumis. & Haute & 4 \\
\hline
US09 & En tant que système, je veux extraire le texte des CV (PyMuPDF sur PDF natif,
OCR Tesseract sur les zones scannées), puis les informations structurées via NER RoBERTa,
persistées dans \texttt{AnalyseCV}. & Haute & 5 \\
\hline
US10 & En tant que système, je veux générer un résumé LLM (Groq), calculer la similarité
sémantique CV/offre (Sentence Transformers), produire un score composite (compétences,
expérience, raisonnement LLM) et classer les candidats (ChromaDB, FAISS), via le service
Python/FastAPI. & Haute & 5 \\
\hline
US11 & En tant que tenant, je veux planifier un quiz technique généré par IA, accessible
via \texttt{QuizToken}, avec résultats (\texttt{QuizResult}) et indicateurs de fraude
consultables. & Haute & 6 \\
\hline
US12 & En tant que tenant, je veux générer un rapport d'entretien IA exportable en PDF et
exporter les données agrégées en Excel et Word. & Moyenne & 6 \\
\hline
US13 & En tant que candidat, je veux passer un entretien virtuel immersif avec un avatar
3D animé, interagir vocalement (TTS/STT) et recevoir une évaluation sémantique automatique. & Moyenne & 7 \\
\hline
US14 & En tant que système, je veux détecter les anomalies comportementales en temps réel
(changement d'onglet, absence caméra, vérification faciale) lors des quiz
et entretiens, et les intégrer aux rapports (\texttt{AlertesSecurite},
\texttt{NbChangementsOnglet}, \texttt{VerificationFacialeOk}). & Haute & 7 \\
\hline
US15 & En tant que tenant, je veux planifier les entretiens virtuels (créneaux,
\texttt{LienToken}) et déployer la plateforme via Docker et CI/CD (GitHub Actions). & Moyenne & 7 \\
\hline
\end{tabular}
\caption{Backlog du produit TalentFlow}
\label{tab:backlog}
\end{table}

% ============================================================
\section{Diagramme de Classes}
% ============================================================

Le diagramme de classes (figure~\ref{fig:classes}) modélise la structure statique de
TalentFlow. Les principales entités sont :

\begin{itemize}
    \item \textbf{Utilisateur :} entité centrale avec rôle (\texttt{RoleUtilisateur} :
    Admin, Tenant, Expert, Candidat), informations d'identité et d'authentification
    externe (\texttt{ExternalId}, \texttt{AuthProvider}).

    \item \textbf{Tenant :} lié à un \texttt{Utilisateur}, porte un statut
    (\texttt{TenantStatus} : Pending, Approved, Rejected) et possède un
    \texttt{TenantProfile} (titre, réseaux sociaux, stack technique, statut de recrutement, logo).

    \item \textbf{Expert :} collaborateur RH d'un tenant et d'une entreprise, avec rôle
    (\texttt{ExpertRole} : Evaluator, SeniorEvaluator, Manager) et statut (\texttt{ExpertStatus}).

    \item \textbf{Entreprise :} société rattachée à un tenant (nom, secteur, RNE, logo).

    \item \textbf{OffreEmploi :} offre publiée par une entreprise avec type de contrat
    (\texttt{TypeContrat} : CDI, CDD, Freelance, Stage, Alternance, Intérim), token public
    (\texttt{PublicToken}) et date limite de candidature.

    \item \textbf{FormulaireCandidature / ChampPersonnalise :} formulaire configurable
    avec champs de types : Texte, Nombre, Date, ChoixMultiple, Fichier.

    \item \textbf{Candidat :} lié à un \texttt{Utilisateur}, possède un
    \texttt{CandidateProfile} (compétences, expériences, formations, certifications,
    langues, CV, avatar, token quiz).

    \item \textbf{SavedJob :} offre sauvegardée en favori par un candidat.

    \item \textbf{Candidature :} soumission pour une offre, portant le CV (\texttt{CvUrl})
    et les réponses au formulaire (\texttt{FormulaireResponses}).

    \item \textbf{AnalyseCV :} résultat IA : texte extrait (\texttt{TexteExtrait}),
    compétences, expériences, certifications, classification (détail du scoring),
    score global (\texttt{Score}) et résumé LLM (\texttt{Resume}).

    \item \textbf{Quiz :} quiz technique avec \texttt{QuizToken}, questions JSON générées
    par IA, minuteur et date de planification.

    \item \textbf{QuizResult :} résultat : score, total, pourcentage, grade, réponses
    (\texttt{AnswersJson}) et temps par question (\texttt{TimeSpentJson}).

    \item \textbf{EntretienIA :} entretien avec créneaux disponibles (\texttt{CreneauxDisponibles}),
    \texttt{LienToken}, score et indicateurs de sécurité : \texttt{AlertesSecurite},
    \texttt{NbChangementsOnglet}, \texttt{DureeMinutes}, \texttt{VerificationFacialeOk}
    et rapport IA final (\texttt{RapportIA}).

    \item \textbf{AvisExpert :} évaluation d'un membre sur une candidature (score et
    commentaire).

    \item \textbf{AssignationExpert :} association Expert — OffreEmploi pour délégation
    de l'évaluation.

    \item \textbf{Notification / CandidatNotification :} notifications destinées aux
    experts et aux candidats.
\end{itemize}

\begin{figure}[ht]
    \centering
    \includegraphics[width=\textwidth]{images/DiagrammeClasses_PFE.png}
    \caption{Diagramme de classes de TalentFlow}
    \label{fig:classes}
\end{figure}

% ============================================================
\section{Environnement de Travail dans ce Stage}
% ============================================================

\subsection{Frameworks et Langages Utilisés}

Le tableau~\ref{tab:frameworks} présente les technologies principales utilisées pour
le développement de TalentFlow.

\begin{table}[h]
\centering
\renewcommand{\arraystretch}{1.5}
\begin{tabular}{|p{3cm}|p{3cm}|p{7cm}|}
\hline
\textbf{Technologie} & \textbf{Version} & \textbf{Rôle} \\
\hline
Vue.js & 3.x & Framework frontend JavaScript (Vite, composition API). \\
\hline
.NET / C\# & .NET 10 & Framework backend pour les APIs RESTful (ASP.NET Core). \\
\hline
Python / FastAPI & 3.11+ & Microservice IA : orchestration des pipelines NLP et exposition REST. \\
\hline
PostgreSQL & 15+ & Base de données relationnelle avec isolation multi-tenant. \\
\hline
PyMuPDF (fitz) & 1.24+ & Extraction du texte natif des fichiers PDF. \\
\hline
Tesseract OCR & 5.x & Reconnaissance optique sur pages ou zones PDF scannées. \\
\hline
spaCy / RoBERTa & 3.x & NER fine-tuné pour l'extraction structurée des CV. \\
\hline
Sentence Transformers & 3.x & Embeddings et similarité sémantique CV $\leftrightarrow$ offre
(\textit{paraphrase-multilingual-MiniLM-L12-v2}). \\
\hline
ChromaDB & 0.5+ & Base vectorielle pour la persistance et la recherche de profils candidats. \\
\hline
FAISS & 1.8+ & Indexation et classement rapide des candidats par offre. \\
\hline
Groq API & --- & LLM pour résumé de profil, raisonnement du score et génération de contenus IA. \\
\hline
Three.js & 0.170+ & Moteur de rendu 3D WebGL pour l'avatar recruteur virtuel. \\
\hline
TensorFlow.js & 4.x & Détection de signaux comportementaux en temps réel côté client. \\
\hline
Docker & 24+ & Conteneurisation de l'ensemble des composants de la plateforme. \\
\hline
\end{tabular}
\caption{Frameworks et langages utilisés}
\label{tab:frameworks}
\end{table}

\subsection{Environnement Logiciel}

\begin{table}[h]
\centering
\renewcommand{\arraystretch}{1.5}
\begin{tabular}{|p{4cm}|p{9cm}|}
\hline
\textbf{Outil} & \textbf{Usage} \\
\hline
Visual Studio 2022 & IDE principal pour le développement de l'API .NET (C\#). \\
\hline
Visual Studio Code & Éditeur pour le frontend Vue.js, le service Python/FastAPI et les fichiers de configuration. \\
\hline
GitHub & Gestion du code source, versioning et collaboration. \\
\hline
GitHub Actions & Pipeline CI/CD pour l'automatisation des builds et des déploiements. \\
\hline
Docker Desktop & Conteneurisation et orchestration locale des services. \\
\hline
Postman & Test et documentation des APIs RESTful. \\
\hline
pgAdmin & Administration et visualisation de la base de données PostgreSQL. \\
\hline
Ready Player Me & Conception des avatars 3D exportés au format \texttt{.glb}. \\
\hline
Mixamo (Adobe) & Bibliothèque d'animations 3D (idle, gestes, transitions) pour l'avatar. \\
\hline
Cloudinary & Hébergement cloud des CV et médias des candidats. \\
\hline
\end{tabular}
\caption{Environnement logiciel}
\label{tab:logiciel}
\end{table}

\subsection{Exemple d'Utilisation}

À titre d'illustration, voici le déroulement d'un cycle complet d'analyse de candidature
sur TalentFlow :

\begin{enumerate}
    \item Un candidat postule à une offre via le lien public tokenisé et uploade son CV
    (stocké sur Cloudinary).
    \item L'API .NET enregistre la candidature et délègue le traitement au service IA Python.
    \item Le service IA extrait d'abord le texte natif du PDF avec \textbf{PyMuPDF}. Si le
    document contient des pages ou zones scannées sans texte sélectionnable, \textbf{Tesseract OCR}
    complète l'extraction sur les images associées.
    \item Le modèle \textbf{NER RoBERTa} (spaCy) structure les informations (compétences,
    expériences, diplômes, certifications, langues).
    \item \textbf{Sentence Transformers} calcule la similarité sémantique entre le profil
    candidat et l'offre ; \textbf{ChromaDB} et \textbf{FAISS} permettent le matching
    et le classement des candidatures pour une même offre.
    \item Le LLM (\textbf{Groq}) génère un résumé professionnel et un score composite
    (compétences, expérience, adéquation sémantique, raisonnement) ; les résultats sont
    persistés dans \texttt{AnalyseCV}.
    \item Le recruteur consulte le score et le résumé depuis son tableau de bord, puis
    planifie un quiz technique ou un entretien virtuel.
    \item Lors du quiz ou de l'entretien, le navigateur surveille les anomalies comportementales
    (changements d'onglet, caméra, vérification faciale) ; les indicateurs sont intégrés au rapport.
    \item Un rapport complet incluant les indicateurs de sécurité est généré et exportable en PDF.
\end{enumerate}

\subsection{Environnement Matériel}

\begin{table}[h]
\centering
\renewcommand{\arraystretch}{1.5}
\begin{tabular}{|p{4cm}|p{9cm}|}
\hline
\textbf{Composant} & \textbf{Caractéristiques} \\
\hline
Processeur & Intel Core i7 (12\textsuperscript{e} génération) ou équivalent \\
\hline
Mémoire RAM & 16 Go DDR4 \\
\hline
Stockage & SSD 512 Go \\
\hline
Système d'exploitation & Windows 11 \\
\hline
Navigateur & Google Chrome / Microsoft Edge (support WebGL et Web Speech API requis) \\
\hline
\end{tabular}
\caption{Environnement matériel}
\label{tab:materiel}
\end{table}

% ============================================================
\section{Architecture Globale}
% ============================================================

\subsection{Architecture Physique}

L'architecture physique de TalentFlow repose sur une conteneurisation Docker de
l'ensemble des composants, déployables sur toute infrastructure cloud :

\begin{itemize}
    \item \textbf{Conteneur Frontend :} application Vue.js 3.x servie via Nginx,
    intégrant Three.js (rendu 3D avatar), TensorFlow.js (détection de signaux
    comportementaux) et la Web Speech API (STT).

    \item \textbf{Conteneur API .NET :} service ASP.NET Core exposant les APIs RESTful,
    orchestrant la logique métier, l'authentification JWT et le filtrage multi-tenant ;
    il délègue les traitements IA au microservice Python.

    \item \textbf{Conteneur Service IA :} microservice Python/FastAPI gérant l'extraction
    PDF (PyMuPDF), l'OCR Tesseract (pages scannées), le NER RoBERTa, les embeddings
    (Sentence Transformers), le matching vectoriel (ChromaDB, FAISS), le LLM (Groq),
    le scoring, la génération de questions, l'évaluation sémantique, la prédiction de
    fraude et la génération des rapports.

    \item \textbf{Conteneur PostgreSQL :} base de données relationnelle.

    \item \textbf{Cloudinary (externe) :} service cloud de stockage des CV et médias.
\end{itemize}

\subsection{Architecture Logique}

L'architecture logique de TalentFlow s'articule autour de quatre couches distinctes :

\begin{itemize}
    \item \textbf{Couche Présentation (Vue.js 3.x) :} interfaces Guest, Candidat, Expert,
    Tenant et Admin. Intègre Three.js pour l'avatar 3D, TensorFlow.js et la vérification
    faciale pour la détection de fraude côté client, ainsi que la Web Speech API (STT).
    Communique avec l'API .NET via JWT.

    \item \textbf{Couche Métier (API .NET) :} orchestre la logique métier, applique
    l'isolation multi-tenant (filtrage par \texttt{TenantId} via les entités liées :
    entreprise, expert, offre) et délègue les traitements intelligents au service Python.

    \item \textbf{Couche IA (Python/FastAPI) :} expose des endpoints REST pour l'extraction
    PyMuPDF, l'OCR Tesseract, le NER RoBERTa, les embeddings Sentence Transformers,
    le matching ChromaDB/FAISS, le résumé et le raisonnement via Groq, le scoring composite,
    la génération de questions, l'évaluation sémantique des réponses, la prédiction de
    fraude et la génération des rapports JSON.

    \item \textbf{Couche Données (PostgreSQL + Cloudinary) :} persistance relationnelle
    avec isolation logique par tenant ; médias hébergés sur Cloudinary.
\end{itemize}

TalentFlow adopte le modèle \textbf{base de données partagée, schéma partagé}, avec
isolation logique par tenant : les entités métier (\texttt{Entreprise}, \texttt{Expert},
\texttt{OffreEmploi}, etc.) sont rattachées à un tenant via des clés étrangères
(\texttt{TenantId} ou chaîne offre $\rightarrow$ entreprise $\rightarrow$ tenant).
L'API .NET applique ce filtrage dans les services et contrôleurs à partir des claims JWT.

% ============================================================
\section{Sprint Backlog Général de Notre Application}
% ============================================================

La planification du projet a été réalisée en \textbf{8 sprints} (Sprint 0 à Sprint 7).
Le tableau~\ref{tab:sprints} résume le découpage et les objectifs de chaque sprint.

\begin{table}[h]
\centering
\renewcommand{\arraystretch}{1.5}
\begin{tabular}{|c|p{9cm}|c|c|}
\hline
\textbf{Sprint} & \textbf{Objectifs} & \textbf{Début} & \textbf{Fin} \\
\hline
Sprint 0 & Apprentissage technologique, design system et maquettes UI/UX, modélisation
UML, schéma de base de données PostgreSQL, conception de l'architecture SaaS multi-tenant. & [Date] & [Date] \\
\hline
Sprint 1 & Infrastructure multi-tenant, gestion des utilisateurs et des rôles
(Admin, Tenant, Expert, Candidat), espace administrateur. & [Date] & [Date] \\
\hline
Sprint 2 & Espace tenant : gestion des membres (\texttt{Expert}), des entreprises
(\texttt{Entreprise}) et des \texttt{TenantProfile}. & [Date] & [Date] \\
\hline
Sprint 3 & Module de gestion des offres : création, personnalisation des
\texttt{FormulaireCandidature} et \texttt{ChampPersonnalise}, publication et
génération de liens publics tokenisés (\texttt{PublicToken}). & [Date] & [Date] \\
\hline
Sprint 4 & Interface publique des offres, dépôt de candidature, upload CV (Cloudinary),
espace candidat (\texttt{CandidateProfile}, \texttt{SavedJob}), suivi du pipeline. & [Date] & [Date] \\
\hline
Sprint 5 & Pipeline IA : extraction PDF (PyMuPDF), OCR Tesseract (zones scannées),
NER RoBERTa, similarité Sentence Transformers, matching ChromaDB/FAISS, résumé et
scoring LLM (Groq), persistance dans \texttt{AnalyseCV}. & [Date] & [Date] \\
\hline
Sprint 6 & Quiz IA (\texttt{Quiz}, \texttt{QuizResult}), assignation experts
(\texttt{AssignationExpert}), avis experts (\texttt{AvisExpert}), génération
de rapports, export PDF / Excel / Word. & [Date] & [Date] \\
\hline
Sprint 7 & Entretien virtuel (avatar 3D, TTS/STT, lip-sync), évaluation sémantique,
détection de fraude pour quiz et entretiens (\texttt{EntretienIA}),
planification, Docker et CI/CD (GitHub Actions). & [Date] & [Date] \\
\hline
\end{tabular}
\caption{Sprint Backlog général de TalentFlow}
\label{tab:sprints}
\end{table}

% TODO : Insérer le diagramme de Gantt.
% \begin{figure}[h]
%     \centering
%     \includegraphics[width=\textwidth]{images/gantt.png}
%     \caption{Diagramme de Gantt du projet}
%     \label{fig:gantt}
% \end{figure}

\section*{Conclusion}

Ce chapitre a permis de formaliser l'ensemble des exigences de TalentFlow. L'identification
des cinq acteurs humains et de l'acteur système AI Service a structuré l'analyse des
besoins fonctionnels. Le diagramme de cas d'utilisation global, le backlog du produit
et le diagramme de classes ont formalisé ces besoins en cohérence avec le modèle de données.
L'environnement de travail détaillé présente les technologies (dont Sentence Transformers,
ChromaDB, FAISS et Groq pour le pipeline de matching et de scoring), outils logiciels
et matériels retenus. L'architecture globale — physique et logique — a posé les fondations techniques
de la plateforme, avec le sprint backlog général structurant le développement en huit
itérations. Les chapitres suivants présenteront la réalisation concrète de TalentFlow,
sprint par sprint.
